\chapter{Definiciones y abreviaturas}
\label{cap:definiciones}

Este apartado recoge la terminología específica, definiciones y abreviaturas empleadas a lo largo del documento, con el fin de facilitar su comprensión.

\begin{table}[h!]
\centering
\small
\begin{tabular}{|p{3.5cm}|p{11cm}|}
\hline
\textbf{Término} & \textbf{Definición} \\ \hline
GTFS & \textit{General Transit Feed Specification}. Estándar abierto para publicar información de transporte público, como rutas, paradas, viajes, horarios y calendarios. \\ \hline
GTFS \textit{Realtime} & Extensión de GTFS orientada a información dinámica del servicio, como retrasos, alertas, cancelaciones o posiciones de vehículos. \\ \hline
API & \textit{Application Programming Interface}. Interfaz que permite a una aplicación acceder a datos o funcionalidades proporcionadas por otro sistema. \\ \hline
REST & \textit{Representational State Transfer}. Estilo arquitectónico para servicios web basado en HTTP. \\ \hline
JWT & \textit{JSON Web Token}. Estándar para la transmisión segura de información entre cliente y servidor en forma de objeto JSON firmado. \\ \hline
CTAN & Red de Consorcios de Transporte de Andalucía, fuente de datos abiertos de transporte público en la comunidad autónoma. \\ \hline
EPSJ & Escuela Politécnica Superior de Jaén. \\ \hline
TFG & Trabajo Fin de Grado. \\ \hline
UJA & Universidad de Jaén. \\ \hline
LTS & \textit{Long Term Support}. Versión de un software con soporte a largo plazo. \\ \hline
ODbL & \textit{Open Database License}, licencia abierta utilizada por OpenStreetMap. \\ \hline
\end{tabular}
\caption{Definiciones y abreviaturas utilizadas en el documento.}
\label{tab:definiciones}
\end{table}
